\documentclass[11pt]{article}
\usepackage[T1]{fontenc}
\usepackage{lmodern}
\usepackage[margin=1in]{geometry}
\usepackage{amsmath,amssymb,amsthm,mathtools}
\usepackage[hidelinks]{hyperref}
\hypersetup{pdftitle={Extremes and Moments in Greedy Coupon Collection},pdfauthor={Christopher D. Long}}
\newtheorem{theorem}{Theorem}[section]
\newtheorem{lemma}[theorem]{Lemma}
\newtheorem{proposition}[theorem]{Proposition}
\newtheorem{corollary}[theorem]{Corollary}
\theoremstyle{remark}\newtheorem{remark}[theorem]{Remark}
\newcommand{\Pbb}{\mathbb P}
\newcommand{\Ebb}{\mathbb E}
\newcommand{\Var}{\operatorname{Var}}
\newcommand{\Cov}{\operatorname{Cov}}
\newcommand{\Fcal}{\mathcal F}
\newcommand{\Rbb}{\mathbb R}
\newcommand{\one}{\mathbf 1}
\newcommand{\dd}{\,\mathrm d}
\newcommand{\norm}[1]{\lVert #1\rVert}
\title{Extremes and Moments in Greedy Coupon Collection}
\author{Christopher D. Long}
\date{}
\begin{document}
\maketitle
\begin{abstract}
A collector is offered a uniformly sampled batch of coupon types and keeps a
least-collected type. For every fixed number of required copies, we prove a
Poisson limit for the centered completion times of individual types, joint
limits for the last completion times, and convergence of all their fixed
absolute and mixed polynomial moments. In particular, we identify the linear
term in the mean and the leading variance. The proof couples the remaining
collection times, after a fixed fluid time, to independent gamma variables;
the common clock error is bounded in every fixed $L^p$ norm by the fraction
of unfinished types. This avoids estimates on a growing fluid-time interval.
The same argument covers sampling with or without replacement, bounded random
batch sizes, and an additional independent holding probability that may tend
to one.
\end{abstract}

\section{Model and results}
There are $n$ types, initially with zero copies. Fix an integer $m\ge1$.
At each trial draw an independent batch size $J$ with a fixed law $\nu$ on
$\{0,\ldots,d\}$, where $d<\infty$ and
\[
 \mu=\Ebb J>0,\qquad p_+=\Pbb(J>0)>0.
\]
Conditional on $J$, sample either a uniform $J$-subset (without replacement)
or $J$ independent uniform types (with replacement). The convention is fixed
throughout the process. A nonempty batch contributes one copy of a
least-collected offered type; an empty batch contributes none. Repeated
offers of one type still contribute at most one copy. Any tie rule using only
present and past information is allowed. All parameters above are fixed as
$n\to\infty$; in the subset model take $n\ge d$.

Let $C_{i,n}$ be the trial when type $i$ first has $m$ copies, and write
\[
 D_{1,n}>D_{2,n}>\cdots>D_{n,n}
\]
for these completion times in decreasing order. In particular, $D_{1,n}$ is
the time to collect $m$ complete sets. Put
\begin{equation}\label{eq:constants}
\begin{gathered}
 c=m-1,\qquad B_n=\log n+c\log\log n,\qquad
 f(x)=\frac{1-\Ebb(1-x)^J}{\mu},\\
 a_\nu=\exp\left\{\int_0^1\left(\frac1{f(x)}-\frac1x\right)\dd x\right\},
 \qquad \lambda=\frac{a_\nu}{c!}.
\end{gathered}
\end{equation}
The integral is finite: $f(x)=x+O(x^2)$ near zero and $f(x)>0$ on $(0,1]$.
All logarithms are natural. Let $E_1,E_2,\ldots$ be independent exponential
variables of rate one and $S_j=E_1+\cdots+E_j$.

\begin{theorem}[Completion extremes and moments]\label{thm:main}
For either sampling convention, the point measure
\begin{equation}\label{eq:ppp}
 \Xi_n=\sum_{i=1}^n\delta_{\mu C_{i,n}/n-B_n}
 \ \Longrightarrow\ \Xi,
 \qquad \text{$\Xi$ Poisson with intensity $\lambda e^{-y}\dd y$,}
\end{equation}
in the vague topology on locally finite point measures on $\Rbb$.
For each fixed $r\ge1$,
\begin{equation}\label{eq:order}
 (\mu D_{j,n}/n-B_n)_{j=1}^r
 \ \Longrightarrow\ (\log\lambda-\log S_j)_{j=1}^r.
\end{equation}
Every fixed absolute moment of the norm of this vector converges, as does
every mixed polynomial moment of its coordinates. Consequently,
\begin{align}
 \Pbb(\mu D_{1,n}/n-B_n\le y)&\longrightarrow
       \exp\{-\lambda e^{-y}\},\label{eq:gumbel}\\
 \Ebb D_{j,n}&=\frac n\mu\bigl(B_n+\log\lambda-\psi(j)+o(1)\bigr),
       \label{eq:mean}\\
 \Cov(D_{j,n},D_{k,n})&=\left(\frac n\mu\right)^2
       \bigl(\psi_1(\max\{j,k\})+o(1)\bigr),\label{eq:cov}
\end{align}
where $\psi$ and $\psi_1$ are the digamma and trigamma functions.
In particular the mean correction for $D_{1,n}$ is
$\log(a_\nu/c!)+\gamma$ and its variance is
$(n/\mu)^2(\pi^2/6+o(1))$.
The consecutive limiting gaps $\log(S_{j+1}/S_j)$, $1\le j<r$, are independent
exponentials of rates $j$; their mixed polynomial moments also arise as limits.
\end{theorem}

\begin{corollary}[Fixed batch size: the Xu--Tang model]\label{cor:xt}
For $J\equiv g\ge1$, $\mu=g$ and
\begin{equation}\label{eq:ag}
 a_g=\exp\{-\gamma-\psi(1/g)-\log g\}.
\end{equation}
Thus Theorem~\ref{thm:main} applies with $\lambda=a_g/(m-1)!$.
For example, $a_1=1$, $a_2=2$, and
$a_3=\sqrt3\exp(\pi/(2\sqrt3))$.
\end{corollary}

Xu and Tang~\cite{XuTang} studied the distinct-batch model and obtained
$\Ebb D_{1,n}=(n/g)(\log n+(m-1)\log\log n)+O(mn)$.
For $g=1$, the limiting law reduces to the classical multiple coupon
collector law of Erd\H{o}s and R\'enyi~\cite{ErdosRenyi}.
Point-process versions for this classical model, including multiple copies,
were proved by Ilienko~\cite{Ilienko}.
Doumas and Spektor~\cite[Theorems 14, 20, 23]{DoumasSpektor} treat the
one-copy greedy case $m=1$, including refined mean asymptotics, the variance,
and a Gumbel limit. Their constant $C_g$ equals $(\log a_g)/g$, so their
centering agrees exactly with ours when $m=1$.
Our result treats every fixed $m$, the full process of late completions and
its moments, and the batch-size and sampling extensions above. We make no
priority claim for the one-copy mean, distribution or variance, or for the
classical $g=1$ point-process limit. Group-drawing Poisson limits such as
those of Schilling and Henze~\cite{SchillingHenze} count all sampled types;
their retention rule differs from the greedy keep-one rule studied here.
Fluid methods for coupon-type hitting problems also appear in
Gast~\cite{Gast}; here the greedy dynamics are interacting, and a comparison
of their clocks supplies the final passage from fixed-time fluid behavior
to extinction. The proof below is self-contained apart from standard
properties of gamma integrals and elementary martingale inequalities,
whose required estimates are given explicitly.

\section{The histogram chain and a fixed-time fluid limit}
Poissonize trials at rate $n/\mu$ and use $h$ for this continuous time.
Loads may be capped at $m$: doing so leaves all counts below $m$, and all
completion times, unchanged. Let $L_j(h)$ count types with load at most $j$,
$0\le j\le c$, and set $L_{-1}=0$.
Define
\begin{equation}\label{eq:Fn}
 F_n(k)=\frac n\mu\begin{cases}
  1-\displaystyle\sum_{\ell=0}^d\nu_\ell
       \dfrac{\binom{n-k}{\ell}}{\binom n\ell},&\text{without replacement},\\[6pt]
  1-\displaystyle\sum_{\ell=0}^d\nu_\ell(1-k/n)^\ell,&\text{with replacement}.
 \end{cases}
\end{equation}
Use $\binom a0=1$, including $a=0$.
A transition that increments a type of load $j<m$ has total rate
\begin{equation}\label{eq:classrate}
 F_n(L_j)-F_n(L_{j-1}).
\end{equation}
Indeed, a batch must avoid all types of load below $j$ and hit at least one
of load $j$. This event, and the resulting histogram, do not depend on the
tie rule. Thus the entire histogram process, including the unlabelled
completion point process, has the same law for every allowed tie rule.
We may henceforth break ties uniformly among distinct types of minimum load.
Every type in a nonempty class of size $N_j=L_j-L_{j-1}$ then has transition rate
\begin{equation}\label{eq:q}
 q_j=\frac{F_n(L_j)-F_n(L_{j-1})}{L_j-L_{j-1}}.
\end{equation}

\begin{lemma}[Rate estimates]\label{lem:rates}
There is a constant $C$ depending only on $d$ such that, if $K=L_c$,
\begin{equation}\label{eq:rates}
 0\le q_j\le1,\qquad q_j\ge1-CK/n
 \quad\text{when }K\le n/2\text{ and }n\ge2d.
\end{equation}
Moreover $F_n(k)\ge\alpha k$, where $\alpha=p_+/\mu>0$, and uniformly in
$0\le k\le n$,
\begin{equation}\label{eq:fluidapprox}
 F_n(k)/n=f(k/n)+O(1/n).
\end{equation}
\end{lemma}
\begin{proof}
For a deterministic nonzero batch size $\ell$ in the subset model,
\[
 F_n(k+1)-F_n(k)
 =\frac\ell\mu\frac{\binom{n-k-1}{\ell-1}}{\binom{n-1}{\ell-1}}
 =\frac\ell\mu\prod_{a=1}^{\ell-1}\left(1-\frac{k}{n-a}\right)
\]
whenever $k\le n-\ell$; the binomial formula remains valid for all $k<n$.
After averaging with weights $\nu_\ell$, these increments are in $[0,1]$.
For $k<K\le n/2$ and $n\ge2d$, the product is at least
$1-2(\ell-1)K/n$. In the replacement model the mean value theorem gives
an increment $(\ell/\mu)(1-\theta/n)^{\ell-1}$ with
$k<\theta<k+1\le K$, which has the same bounds. Since $q_j$ is an average
of increments, one may take $C=2d$.
A nonempty batch hits a set of $k$ types with probability at least $k/n$,
proving the global lower bound. Finally, the fixed-length product defining
the hypergeometric avoidance probability differs uniformly from
$(1-k/n)^\ell$ by $O(1/n)$. For replacement the approximation is exact.
\end{proof}

Let $v_{-1}=0$ and define the deterministic profile by
\begin{equation}\label{eq:ode}
 v_j'=-\{f(v_j)-f(v_{j-1})\},\qquad v_j(0)=1,
 \qquad 0\le j\le c.
\end{equation}

\begin{lemma}[Fixed-time concentration and eventual sparsity]\label{lem:fluid}
For every fixed $b<\infty$ and $\varepsilon>0$, there are positive constants
$A_{b,\varepsilon},a_{b,\varepsilon}$ such that, for all sufficiently large $n$,
\begin{equation}\label{eq:concentration}
 \Pbb\left(\max_{0\le j\le c}\sup_{0\le h\le b}
       |L_j(h)/n-v_j(h)|>\varepsilon\right)
 \le A_{b,\varepsilon}e^{-a_{b,\varepsilon}n}.
\end{equation}
The solution satisfies $0<v_0(h)\le\cdots\le v_c(h)\le1$, every $v_j$
is nonincreasing and tends to zero, and
\begin{equation}\label{eq:v0}
 e^h v_0(h)\longrightarrow a_\nu.
\end{equation}
\end{lemma}
\begin{proof}
The polynomial vector field in \eqref{eq:ode} is Lipschitz on the compact
ordered simplex $0\le v_0\le\cdots\le v_c\le1$ and points inward on its
boundary. For example, at $v_j=v_{j-1}$ the derivative of their difference
is $f(v_{j-1})-f(v_{j-2})\ge0$. Hence the solution exists globally in that
simplex. Since $f$ is strictly increasing on $[0,1]$, each derivative is
nonpositive. The limits of successive coordinates must agree: if
$v_j(\infty)>v_{j-1}(\infty)$, its derivative stays strictly negative.
Starting with $v_{-1}=0$ proves that every limit is zero. Also $f(x)\le x$
and $v_0'=-f(v_0)$, so $v_0(h)\ge e^{-h}>0$.

For concentration, write each normalized count as its drift integral plus
a martingale $M_j$. The jumps of $M_j$ have absolute size $1/n$ and its
predictable quadratic variation up to $b$ is at most $b/(\mu n)$, since the
trial rate is $n/\mu$. The compensated counting-process exponential
supermartingale and $e^z-1-z\le z^2/[2(1-|z|/3)]$ for $|z|<3$ give
\[
 \Pbb\left(\sup_{h\le b}|M_j(h)|>z\right)
 \le2\exp\left\{-\frac{nz^2}{2(b/\mu+z/3)}\right\}.
\]
For completeness, apply the exponential supermartingale with parameter
$\theta n$ up to the first crossing of $z$, bound its compensator by
$n\theta^2b/[2\mu(1-\theta/3)]$, and optimize at
$\theta=z/(b/\mu+z/3)$; do the same for $-M_j$.
Equation~\eqref{eq:fluidapprox}, a union bound over the $m$ coordinates,
and Gronwall's inequality prove \eqref{eq:concentration}.

Finally, separation of variables in $v_0'=-f(v_0)$ gives
\[
 h=\int_{v_0(h)}^1\frac{\dd x}{f(x)},\qquad
 h+\log v_0(h)=\int_{v_0(h)}^1
       \left(\frac1{f(x)}-\frac1x\right)\dd x.
\]
Taking the limit proves \eqref{eq:v0}.
\end{proof}

\begin{proposition}[All deterministic lower-tail profiles]\label{prop:profiles}
For every fixed $j\ge0$, the solution of \eqref{eq:ode} satisfies
\[
 v_j(h)\sim a_\nu\frac{h^j}{j!}e^{-h}\qquad(h\to\infty).
\]
\end{proposition}
\begin{proof}
Fix $\rho\in(1/2,1)$. Since $v_j(h)\to0$ and $f(x)=x+O(x^2)$,
for all sufficiently large $h$,
$v_j'\le-\rho v_j+v_{j-1}$. Induction and variation of constants give
$v_j(h)=O((1+h)^j e^{-\rho h})$ for each of the finitely many layers in
question. Expanding the equation more accurately gives, for $j\ge1$,
\[
 (e^h v_j(h))'=e^h v_{j-1}(h)
       +O\bigl(e^h(v_j(h)^2+v_{j-1}(h)^2)\bigr).
\]
The error is integrable because $2\rho>1$. Starting with
\eqref{eq:v0}, integration and induction therefore give
$e^h v_j(h)\sim a_\nu h^j/j!$.
\end{proof}

\section{A simultaneous comparison with independent gamma clocks}
The following is the main estimate. It concerns all remaining completion
times, not merely their maximum, and supplies its own moment bounds.

\begin{lemma}[Clock comparison]\label{lem:clock}
Assume $n\ge2d$ and condition on the state at a deterministic time $b$.
Suppose exactly $K$
types are unfinished and $K\le\delta n$, where
$\delta\le\min\{1/2,1/(2C)\}$ and $C$ is from Lemma~\ref{lem:rates}.
If type $i$ needs $r_i\in\{1,\ldots,m\}$ more copies, one can construct
conditionally independent $G_i\sim\operatorname{Gamma}(r_i,1)$ and the
actual residual completion times $R_i$ so that
\begin{equation}\label{eq:clockcoupling}
 G_i\le R_i\le G_i+\Delta\quad\text{for every unfinished }i,
 \qquad 0\le\Delta\le\frac{2C}{n}\sum_{i=1}^K G_i.
\end{equation}
For every fixed $p\ge1$,
\begin{equation}\label{eq:clockmoment}
 \bigl(\Ebb[\Delta^p\mid\Fcal_b]\bigr)^{1/p}
 \le C_{m,d,p}K/n.
\end{equation}
For $K=0$, set $\Delta=0$.
\end{lemma}
\begin{proof}
Starting at $b$, use independent unit-rate Poisson processes $Y_i$ for the
unfinished types. When a type is in class $j$, run its internal clock at
rate $q_j$. The random time change
$Y_i(\int_b^{b+t}q_i(h)\dd h)$ constructs the desired finite-state Markov
chain: its coordinate transition rates are exactly \eqref{eq:q}.
After type $i$ finishes, set its virtual clock rate to one. The $r_i$th
arrival time of $Y_i$ is $G_i$; these are independent gamma variables
conditional on the starting state. Continuity of the internal clocks gives
\[
 \int_b^{b+R_i}q_i(h)\dd h=G_i.
\]
Let $K(h)=L_c(h)$. It is nonincreasing, and \eqref{eq:rates} gives
$1/2\le q_i(h)\le1$ while a type is unfinished. Consequently
$G_i\le R_i\le2G_i$, so all completion times are finite. Define
\[
 \Delta=\frac Cn\int_b^\infty K(h)\dd h.
\]
For each type, its total lost clock time is at most $\Delta$, whence
$R_i-G_i\le\Delta$. Moreover, by summing the unfinished indicators,
\[
 \int_b^\infty K(h)\dd h=\sum_{i=1}^K R_i
 \le2\sum_{i=1}^K G_i.
\]
This proves \eqref{eq:clockcoupling}. Minkowski's inequality and
$\norm{\operatorname{Gamma}(r,1)}_p\le
\norm{\operatorname{Gamma}(m,1)}_p$ for $r\le m$ prove
\eqref{eq:clockmoment}. Independence is used to identify the ideal clocks;
the error estimate itself only needs their marginal moment bounds.
\end{proof}

\begin{remark}
Keeping only the constant lower rate $1-C\delta$ would introduce an error
of order $\delta\log n$ at the maximum. The integrated unfinished population
in \eqref{eq:clockcoupling} instead gives an $L^p$ error of order $\delta$,
uniformly in $n$. This is what allows a fixed fluid time followed by the
limits $n\to\infty$ and then $b\to\infty$.
\end{remark}

\begin{lemma}[Coarse global moments]\label{lem:coarse}
Let $T_n$ be the full completion time in Poissonized time and $D_{1,n}$ the
number of trials. For every fixed $p\ge1$,
\begin{equation}\label{eq:coarse}
 \Ebb T_n^p=O((\log n)^p),\qquad
 \Ebb D_{1,n}^p=O((n\log n)^p).
\end{equation}
The same bounds hold uniformly over initial states with nonnegative loads.
\end{lemma}
\begin{proof}
Divide collection into at most $m$ phases, the $j$th ending when every type
has at least $j$ copies. During a phase, if $k$ types remain at the current
minimum load, this number decreases at continuous rate $F_n(k)\ge\alpha k$.
Its phase duration is therefore dominated by
$\sum_{k=1}^n\operatorname{Exp}(\alpha k)$ with independent summands.
This sum has the distribution of the maximum of $n$ independent
$\operatorname{Exp}(\alpha)$ variables, whose tail above
$\alpha^{-1}(\log n+x)$ is at most $e^{-x}$. Conditional domination in each
phase and the $L^p$ triangle inequality prove the first assertion.

In discrete time the success probability is at least $p_+k/n$.
A phase is dominated by a classical coupon collector that is activated with
probability $p_+$ at each trial. After $t$ trials its probability of being
unfinished is at most $n(1-p_+/n)^t\le ne^{-p_+t/n}$.
Integrating this tail and summing the $m$ phases proves the second bound.
Larger starting loads only remove phases or reduce the numbers still at
the current minimum; the bounds remain valid.
\end{proof}

\section{Independent gamma extremes, including moments}
We record the elementary extreme-value calculation with the moment
estimates needed later. It does not assume independence of the initial
histogram.

\begin{lemma}[Mixed gamma arrays]\label{lem:gamma}
Fix $b\ge0$. Given a random initial configuration, let $G_i$ be independent
conditional on it, with integer shapes $r_i\in\{0,1,\ldots,m\}$ and rate one;
shape zero means $G_i=0$. Suppose the number $N_m$ of shape-$m$ variables
satisfies $N_m/n\to v>0$ in probability. Then
\[
 \sum_{i=1}^n\delta_{b+G_i-B_n}
 \Longrightarrow \operatorname{PPP}(\lambda_b e^{-y}\dd y),
 \qquad \lambda_b=e^b v/c!,
\]
and the largest fixed $r$ values converge jointly to
$(\log\lambda_b-\log S_j)_{j=1}^r$.
If, except on an event of probability $O(e^{-a n})$ for some $a>0$,
$N_m\ge\kappa n$ with $\kappa>0$, then these joint limits also hold with
all fixed absolute and mixed polynomial moments.
\end{lemma}
\begin{proof}
For integer $r\ge1$ and $t\ge0$,
\begin{equation}\label{eq:gammatail}
 Q_r(t):=\Pbb(\operatorname{Gamma}(r,1)>t)
       =e^{-t}\sum_{k=0}^{r-1}\frac{t^k}{k!}.
\end{equation}
Thus, for fixed $y$,
\[
 n Q_m(B_n-b+y)\longrightarrow e^{b-y}/c!,\qquad
 n Q_r(B_n-b+y)\longrightarrow0\quad(r<m).
\]
The density gives the analogous uniform convergence on compact $y$-intervals.
For any nonnegative continuous compactly supported $\phi$, condition on the
shapes and write the Laplace functional as
\[
 \prod_{i=1}^n(1-z_{i,n}),\qquad
 z_{i,n}=\Ebb[1-e^{-\phi(b+G_i-B_n)}\mid(r_i)].
\]
Here $\max_i z_{i,n}=O(1/n)$ and
\[
 \sum_i z_{i,n}\ \longrightarrow\
 \lambda_b\int_{\Rbb}(1-e^{-\phi(y)})e^{-y}\dd y
\]
in probability. Expanding $\log(1-z)$, then using boundedness of the
conditional Laplace functional, proves the point-process assertion.
The same product calculation for finitely many disjoint intervals, including
$(y,\infty)$, gives independent Poisson counts. Their means identify the joint
upper order statistics: under $z=\lambda_b e^{-y}$ the limit becomes a
unit-rate Poisson process on $(0,\infty)$, with successive points $S_j$.

Here are explicit tail bounds for moments. Set $t_n=B_n-b$ and let $Z_{j,n}$
be the $j$th largest of $b+G_i-B_n$. For large $n$ and $x\ge0$, a union bound
and \eqref{eq:gammatail} give
\begin{equation}\label{eq:right}
 \Pbb(Z_{1,n}>x)\le C_b(1+x)^c e^{-x}.
\end{equation}
On a shape configuration with $N_m\ge\kappa n$, the number of shape-$m$
variables exceeding $t_n-x$ is binomial with mean at least
$A_b e^x$ for $0\le x\le t_n/2$; this follows from the highest term in
\eqref{eq:gammatail} and $t_n\sim\log n$.
For $X\sim\operatorname{Bin}(N,q)$,
\[
 \Pbb(X<j)\le e^{j-1}\Ebb e^{-X}
       \le \exp\{j-1-(1-e^{-1})Nq\}.
\]
Consequently, uniformly over these shape configurations,
\begin{equation}\label{eq:left}
 \Pbb(Z_{j,n}<-x\mid(r_i))\le
 \begin{cases}
 \exp\{j-1-A'_b e^x\},&0\le x\le t_n/2,\\
 \exp\{j-1-A'_b e^{t_n/2}\},&t_n/2<x\le t_n,\\
 0,&x>t_n.
 \end{cases}
\end{equation}
For the middle interval, use monotonicity and the estimate at $t_n/2$.
The contribution of that interval to every fixed tail integral vanishes;
indeed $t_n^p\exp(-A'_b e^{t_n/2})\to0$.
The first and right tails give uniform integrability of every fixed power.
On the exceptional shape event, negative parts are at most $t_n=O(\log n)$;
positive parts have uniformly bounded conditional moments by
\eqref{eq:right}. Its contribution therefore vanishes.
Weak convergence and these bounds give the stated moments, using a higher
moment and H\"older's inequality for mixed polynomials.
\end{proof}

\section{Passage to the interacting process}
Let $T_{i,n}$ be type $i$'s Poissonized completion time, and let
$T_{1,n}^{\downarrow}>\cdots>T_{n,n}^{\downarrow}$ be their decreasing order.
Fix a sufficiently large but finite $b$, and put
\[
 \delta_b=2v_c(b),\qquad
 A_{n,b}=\{L_c(b)\le\delta_b n,\ L_0(b)\ge v_0(b)n/2\}.
\]
Choose $b$ so that $\delta_b\le\min\{1/2,1/(2C)\}$.
Lemma~\ref{lem:fluid} implies
\begin{equation}\label{eq:good}
 \Pbb(A_{n,b}^c)\le A_b e^{-a_b n},\qquad
 L_0(b)/n\longrightarrow v_0(b)>0.
\end{equation}
Given $\Fcal_b$, construct the gamma variables of Lemma~\ref{lem:clock} for
all unfinished types, even on $A_{n,b}^c$ (the small-error bound is asserted
only on $A_{n,b}$). Give a completed type the ideal time $b$, and an
unfinished type the ideal time $\widehat T_{i,n}=b+G_i$.
Let $\widehat T_{j,n}^{\downarrow}$ denote their decreasing order.

The number of shape-$m$ variables is exactly $L_0(b)$. By
Lemma~\ref{lem:gamma}, the ideal point process and its largest $r$ points
have the stated limits with
\begin{equation}\label{eq:lambdab}
 \lambda_b=\frac{e^b v_0(b)}{c!}\longrightarrow\lambda,
\end{equation}
including all moments. On $A_{n,b}$ and for all sufficiently large $n$ there
are at least $r$ unfinished types. All their actual and ideal times exceed
$b$, while all completed types have actual times at most $b$.
The coordinatewise comparison \eqref{eq:clockcoupling} therefore gives
\begin{equation}\label{eq:sorterror}
 \max_{1\le j\le r}|T_{j,n}^{\downarrow}
                 -\widehat T_{j,n}^{\downarrow}|\le\Delta
       \quad\text{on }A_{n,b}.
\end{equation}
On its complement, each actual order statistic is bounded by $T_n$, and
each ideal one by $b+\max_{i\le n}\operatorname{Gamma}(m,1)$ under an
appropriate conditional coupling. These variables have $2p$th moments
$O_b((\log n)^{2p})$ by Lemma~\ref{lem:coarse} and
\eqref{eq:gammatail}. Cauchy--Schwarz and \eqref{eq:good} show that their
contribution on $A_{n,b}^c$ to a $p$th moment tends to zero.
Combining this fact with \eqref{eq:clockmoment} proves
\begin{equation}\label{eq:Lpapprox}
 \limsup_{n\to\infty}
 \norm{\max_{1\le j\le r}|T_{j,n}^{\downarrow}
                 -\widehat T_{j,n}^{\downarrow}|}_p
 \le C_{m,d,p}\delta_b\longrightarrow0\qquad(b\to\infty).
\end{equation}

For clarity, this gives both weak and moment convergence without assuming
uniform integrability of a quantity known only to be small in probability.
For a bounded Lipschitz function of the largest $r$ points, its expectations
for the two vectors differ by at most a constant times the $p=1$ bound in
\eqref{eq:Lpapprox}. Take $n\to\infty$, then $b\to\infty$, using
\eqref{eq:lambdab}. This proves joint weak convergence. For any fixed power,
choose one larger exponent in \eqref{eq:Lpapprox} and use the ideal moment
convergence. Alternatively, the inequality
\[
 |\norm{x}^p-\norm{y}^p|
 \le p\norm{x-y}(\norm{x}+\norm{y})^{p-1}\quad(p\ge1)
\]
and H\"older's inequality give the moment comparison directly. The limiting
ideal norms stay bounded as $b\to\infty$, since $\log\lambda_b\to\log\lambda$.
Mixed polynomials follow by the same argument; powers between zero and one
follow from a higher moment.

To also obtain the point process, fix a compact interval $I$ and a
continuous test function supported in $I$. On $\{\Delta\le\varepsilon\}$,
all actual or ideal points that can contribute lie among the ideal points
in the $\varepsilon$-enlargement of $I$, and each contributing pair is at
distance at most $\varepsilon$. Completed-before-$b$ points are outside $I$
for large $n$. The ideal counts in the enlargement are tight, by
Lemma~\ref{lem:gamma}. Markov's inequality and \eqref{eq:clockmoment} give
\[
 \limsup_n\Pbb(A_{n,b}\cap\{\Delta>\varepsilon\})
 \le C\delta_b/\varepsilon.
\]
First choose small $\varepsilon$, then large fixed $b$ with
$\delta_b/\varepsilon$ small, and finally let $n\to\infty$.
Uniform continuity of the test function proves convergence of Laplace
functionals to those of intensity $\lambda e^{-y}\dd y$.
This proves the Poissonized version of \eqref{eq:ppp}.

\section{De-Poissonization and evaluation of the limits}
Take the actual discrete trial sequence independently of rate-one
exponential interarrival variables $U_1,U_2,\ldots$. Poissonization means
\[
 T_{i,n}=\frac\mu n\sum_{k=1}^{C_{i,n}}U_k.
\]
Conditional on $D_{1,n}=N$, the centered partial sums of the $U_k$ are a
martingale. For any integer $q\ge1$, the elementary moment expansion of
independent centered summands, followed by Doob's $L^{2q}$ inequality, gives
\[
 \Ebb\max_{k\le N}\left|\sum_{\ell=1}^k(U_\ell-1)\right|^{2q}
 \le C_q N^q\quad(N\ge1).
\]
In the expansion, any nonzero term has no singleton index, so at most $q$
distinct indices occur; their moments are finite. This proves the terminal
moment bound used by Doob. Independence of the trial sequence and the clocks,
and \eqref{eq:coarse}, now imply
\begin{equation}\label{eq:depois}
 \Ebb\max_{k\le D_{1,n}}
 \left|\frac\mu n\sum_{\ell=1}^k(U_\ell-1)\right|^{2q}
 =O\bigl((\log n/n)^q\bigr)\longrightarrow0.
\end{equation}
Every completion time through the last one is changed by at most this
maximum. Thus the joint laws and moments pass to discrete time. Compactly
supported point-process tests pass by the same tight-count and
uniform-continuity argument used above.

The gamma integral gives
$\Ebb\log S_j=\psi(j)$ and $\Var(\log S_j)=\psi_1(j)$.
For $j<k$, the beta--gamma change of variables shows that $S_j/S_k$ is
independent of $S_k$. Hence
$\Cov(\log S_j,\log S_k)=\Var(\log S_k)=\psi_1(k)$,
which proves \eqref{eq:mean}--\eqref{eq:cov}.
For the spacings, the joint density of
$0<s_1<\cdots<s_r$ is $e^{-s_r}$. Changing variables to
$u_j=s_j/s_{j+1}$, $1\le j<r$, and $s_r$ gives the product density of
independent $u_j\sim\operatorname{Beta}(j,1)$ and
$s_r\sim\operatorname{Gamma}(r,1)$. Therefore $-\log u_j$ are independent
exponentials of rates $j$. This completes the proof of
Theorem~\ref{thm:main}.

To verify Corollary~\ref{cor:xt}, for $J\equiv g$ substitute $t=1-x$ in
\eqref{eq:constants}:
\[
 \log a_g=\int_0^1\left(\frac g{1-t^g}-\frac1{1-t}\right)\dd t.
\]
With an upper cutoff $1-\varepsilon$, set $z=t^g$ in the first integral and
separate $z^{1/g-1}-1$. The singular pieces give
$-\log(1-(1-\varepsilon)^g)+\log\varepsilon\to-\log g$.
The remaining integral is $-\psi(1/g)-\gamma$ by the digamma integral
\cite[\S5.9.16]{DLMF}. The displayed special values follow from the same
integral, or the standard digamma values in \cite[\S5.4]{DLMF}.

\section{Further consequences and scope}
\begin{corollary}[A finite sparse-state mean bound]\label{cor:finite}
Start from any deterministic state with $K$ unfinished types, residual quotas
$r_1,\ldots,r_K\le m$, and $K/n\le\min\{1/2,1/(2C)\}$, $n\ge2d$.
Write $D_{\rm rem}$ for its remaining number of trials and
\[
 H(r_1,\ldots,r_K)=\int_0^\infty
   \left[1-\prod_{i=1}^K(1-Q_{r_i}(t))\right]\dd t,
\]
where $Q_r$ is given by \eqref{eq:gammatail}. Then
\[
 0\le\Ebb D_{\rm rem}-\frac n\mu H(r_1,\ldots,r_K)
 \le\frac{2C}{\mu}\sum_{i=1}^K r_i.
\]
For $K=0$ all terms vanish.
\end{corollary}
\begin{proof}
The integral is the mean maximum of the independent residual gamma clocks.
Taking maxima and expectations in \eqref{eq:clockcoupling} bounds the
continuous residual mean between $H$ and
$H+(2C/n)\sum_i r_i$. Independent exponential trial intervals give
$\Ebb T_{\rm rem}=(\mu/n)\Ebb D_{\rm rem}$; finiteness follows from
Lemma~\ref{lem:coarse}. Rescaling proves the claim.
\end{proof}

This bound can close a truncated finite-state computation with a certified
additive error. It does not reduce the state count of the exact histogram
dynamic program of Xu and Tang~\cite{XuTang}. For small parameters, that
recursion remains an exact way to compute the mean and, by a second-moment
recursion, the variance.

\begin{corollary}[The number still unfinished]\label{cor:counts}
At trial $\lfloor(n/\mu)(B_n+y)\rfloor$, the number of types with fewer
than $m$ copies converges in distribution to
$\operatorname{Poisson}(\lambda e^{-y})$.
For fixed $j$,
\[
 \Pbb(\mu D_{j,n}/n-B_n\le y)\longrightarrow
 e^{-\lambda e^{-y}}\sum_{k=0}^{j-1}\frac{(\lambda e^{-y})^k}{k!}.
\]
\end{corollary}
\begin{proof}
Vague convergence alone would not control an unbounded interval.
Instead, truncate the interval $(y,\infty)$ at $M$. The probability that a
point lies above $M$ is the upper tail of the largest point; convergence of
that maximum makes its limiting upper bound tend to zero as $M\to\infty$.
The Poisson limit has no atom at $y$ or $M$. Apply the bounded-interval
point-process limit, then remove the truncation. The displayed formula is
the probability that this Poisson count is at most $j-1$.
The rounding error is at most $\mu/n$ and does not affect these limits.
\end{proof}

\begin{corollary}[Additional independent holding]\label{cor:holding}
Activate each trial independently with probability $p_n\in(0,1]$, performing
one trial of the model above only on activation. Allow $p_n$ to tend to zero.
If $D^{(p)}_{j,n}$ and $C^{(p)}_{i,n}$ count total trials, every assertion
above holds with $\mu D_{j,n}/n$ and $\mu C_{i,n}/n$ replaced by
$p_n\mu D^{(p)}_{j,n}/n$ and $p_n\mu C^{(p)}_{i,n}/n$.
\end{corollary}
\begin{proof}
Let $V_\ell$ be independent geometric variables with success probability
$p_n$ and support $\{1,2,\ldots\}$, independent of the active trial sequence.
Then $D^{(p)}_{j,n}=\sum_{\ell=1}^{D_{j,n}}V_\ell$ and similarly for each
type. The variables $p_nV_\ell-1$ are centered and their absolute moments
of every fixed order are bounded uniformly in $p_n$: for example
$\Pbb(p_nV_\ell>x)\le e^{1-x}$ for $x\ge0$.
The same moment expansion and maximal inequality as in
\eqref{eq:depois} show that the uniform scaled clock error tends to zero in
every fixed $L^p$. Transfer all conclusions as before.
\end{proof}

The proofs require fixed $m$ and fixed bounded positive-mean batch law.
They do not establish growing-$m$ or growing-batch regimes, nonuniform type
weights, a central limit theorem for the moving histogram, or a quantitative
rate for the final Gumbel approximation. A fixed-time concentration theorem
and the integrated clock comparison suffice for the completion results;
no moving-tail normal approximation is asserted or required.

\begin{thebibliography}{9}
\bibitem{XuTang}
W. Xu and A. K. Tang,
A generalized coupon collector problem,
\emph{Journal of Applied Probability} \textbf{48} (2011), 1081--1094.
\href{https://doi.org/10.1239/jap/1324046020}{doi:10.1239/jap/1324046020}.

\bibitem{ErdosRenyi}
P. Erd\H{o}s and A. R\'enyi,
On a classical problem of probability theory,
\emph{Magyar Tud. Akad. Mat. Kutat\'o Int. K\"ozl.} \textbf{6} (1961), 215--220.
\href{https://static.renyi.hu/renyi_cikkek/1961_on_a_classical_problem_of_probability_theory.pdf}{Author archive}.

\bibitem{DoumasSpektor}
A. V. Doumas and S. Spektor,
Two Multi--Draw Coupon Collector models with different retention rules,
preprint (2026).
\href{https://arxiv.org/abs/2607.01463}{arXiv:2607.01463}.

\bibitem{Ilienko}
A. Ilienko,
Convergence of point processes associated with coupon collector's and
Dixie cup problems,
\emph{Electronic Communications in Probability} \textbf{24} (2019), paper 51.
\href{https://doi.org/10.1214/19-ECP263}{doi:10.1214/19-ECP263}.

\bibitem{SchillingHenze}
J. Schilling and N. Henze,
Two Poisson limit theorems for the coupon collector's problem with group drawings,
\emph{Journal of Applied Probability} \textbf{58} (2021), 966--977.
\href{https://doi.org/10.1017/jpr.2021.15}{doi:10.1017/jpr.2021.15}.

\bibitem{Gast}
N. Gast,
Computing hitting times via fluid approximation: application to the coupon
collector problem, preprint (2011).
\href{https://arxiv.org/abs/1107.3385}{arXiv:1107.3385}.

\bibitem{DLMF}
NIST Digital Library of Mathematical Functions,
\href{https://dlmf.nist.gov/5.9}{\S5.9, Integral Representations}, and
\href{https://dlmf.nist.gov/5.4}{\S5.4, Special Values and Extrema}.
\end{thebibliography}
\end{document}
